\chapter{Green Revolution: Gender, Bondage and Agrarian Class}

\section{Impact on Women: Agarwal, Walia and Chaudhary}

\begin{KeyConcept}
\begin{itemize}
\item Agriculture was a \textbf{household activity} (family labour + \textbf{attached labour} --- a family attached to a landlord for life; like the \textbf{gold-mine workers} who were given an advance they could never repay and so stayed attached for life). Green Revolution changed this and created new forms of bondage.
\item \textbf{Bina Agarwal}: technology has a \textbf{clear bias against women} --- women are not given the skills to use the new machines (which are male domains). Agriculture once let women work outside the domestic sphere, but \textbf{mechanisation excluded them} from agrarian work. However, in \textbf{small and marginal farms} (which cannot afford machines) women's role \textbf{increased} --- because they work at a lower wage.
\item The impact on labour: big landlords stopped using \textbf{family labour} (except at peak periods); \textbf{permanent labourers} were retained (for peaks); \textbf{casual labourers} lost work and declined.
\item \textbf{Saila Walla} (Haryana): women labourers are used only in peak periods and for new crops (e.g. picking cotton buds --- a banal activity). Technological change \textbf{strengthened patriarchy} --- the \textbf{khaps} became powerful, control over women's bodies increased, and \textbf{female feticide/infanticide} rose (they got money but not modernity).
\item \textbf{Prem Chaudhary} (khaps, Haryana/Punjab): women had \textbf{no land ownership}, so no access to credit/inputs (no collateral), and were \textbf{neglected by government officials} (who could not believe women can be farmers) and by research/extension services. The answer is \textbf{land ownership rights} --- land is a productive asset giving subsistence, freedom and empowerment.
\end{itemize}
\end{KeyConcept}

\begin{QuickRevision}
\begin{itemize}
\item Bina Agarwal: technology biased against women; mechanisation excluded women; but small/marginal farms $\rightarrow$ women's work increased (low wage).
\item Saila Walla (Haryana): women used in peaks + cotton; patriarchy strengthened (khaps, feticide/infanticide).
\item Prem Chaudhary: no land ownership $\rightarrow$ no credit/inputs; neglected by officials; solution = women's land ownership.
\end{itemize}
\end{QuickRevision}

\section{Feminisation, Domestication and Selective Modernity}

\begin{KeyConcept}
\begin{itemize}
\item \textbf{Feminisation of agriculture}: as men migrated (from non-GR areas to GR areas, and from GR areas to cities), the onus of cultivation fell on women. But this is \textbf{not empowerment} --- women were given the \textbf{least productive sector} of the economy, so it is itself patriarchal.
\item \textbf{Domestication and objectification}: women and luxury cars both became objects of status in GR areas (see Punjabi songs). The \textbf{selective modernity} of GR areas --- they had wealth (Weber's 'new rich'), bought luxury goods and travelled, but remained attached to land and tradition: they accepted TV/cinema but not women going to theatres. After 1991, land acquired at high prices bought Mercedes/BMWs/hotels --- and later debt.
\item \textbf{Violence against women increased} --- with wealth and political power came a sense of \textbf{impunity} (the belief that one will not be convicted; e.g. the lynching video culture).
\item \textbf{Prosperity $\rightarrow$ freedom demands}: the \textbf{Khalistan} movement became powerful after Green Revolution (prosperity, not poverty, fuels demands for freedom --- as in Kashmir).
\item \textbf{Vandana Shiva}: environment and women are linked --- Green Revolution's \textbf{unsustainable agriculture} (monocropping, excessive fertiliser and water) will destroy the environment and further lower women's status; she even argued against using machinery.
\end{itemize}
\end{KeyConcept}

\begin{QuickRevision}
\begin{itemize}
\item Feminisation of agriculture = patriarchal (least productive sector).
\item Selective modernity (consumerism; new rich; khaps; violence; impunity); Khalistan rose with prosperity.
\item Vandana Shiva: unsustainable GR (monocropping/fertiliser/water) $\rightarrow$ environment + women harmed.
\end{itemize}
\end{QuickRevision}

\section{The Rise of New Political Forces and the New Farmers Movement}

\begin{KeyConcept}
\begin{itemize}
\item The dominant castes (Jats, Jat Sikhs, etc.) --- earlier the \textbf{vote bank} behind Congress (a \textbf{patron--client relationship}) --- used their new wealth to become \textbf{leaders themselves}, and mobilised \textbf{on the basis of caste} (in the pursuit of fighting caste, they \textbf{substantialised} it --- Beteille). This rise of the OBCs is called the \textbf{'Silent Revolution'} (Christopher Jaffrelot) --- a smooth, democratic, non-violent transition of power (unlike the violent transition in South India).
\item \textbf{New Farmers Movement} (from the \textbf{New Social Movement} of the 1960s): \textbf{apolitical} (not led by any political party), making \textbf{abstract demands} (not specific, negotiable ones), and raising \textbf{social-reform issues}. Its philosophical basis came from \textbf{Sharad Joshi} (Shetkari Sangathan, Maharashtra) --- a retired civil servant who argued \textbf{'Bharat vs India'} --- the \textbf{internal colonisation} of rural India by urban India (farmers pay market price for inputs but must sell at MSP; they finance urbanisation by cheap food) --- so he demanded a \textbf{free market in agriculture}. Other movements (Tikait in UP; Karnataka's Marxist movement) had Gandhian/Marxist orientations.
\item \textbf{Critique}: \textbf{Gail Omvedt} --- the movement is of the \textbf{enterprising/capitalist farmers}, with a \textbf{class bias} (it ignores agricultural labourers and women). \textbf{Neera Chandhoke} --- it is \textbf{not apolitical} (Tikait supported the BJP, Sharad Joshi backed anti-Congress candidates). \textbf{Mridula Mukherjee} --- there is \textbf{nothing new}: only the leaders changed from peasants to big farmers.
\end{itemize}
\end{KeyConcept}

\begin{QuickRevision}
\begin{itemize}
\item OBC rise = 'Silent Revolution' (Jaffrelot); patron-client/vote-bank $\rightarrow$ caste mobilisation (substantialisation --- Beteille).
\item New Farmers Movement: apolitical, abstract demands, social reform; Sharad Joshi (Shetkari Sangathan, 'Bharat vs India', internal colonisation, free market).
\item Critique: Gail Omvedt (class bias), Neera Chandhoke (not apolitical), Mridula Mukherjee (nothing new).
\end{itemize}
\end{QuickRevision}

\section{The Jajmani System and Its Disintegration}

\begin{KeyConcept}
\begin{itemize}
\item The \textbf{jajmani system} rests on four terms: \textbf{Jajman} (the one receiving service), \textbf{Purohit} (rituals), \textbf{Mahabrahmin} (death rituals, for 10 days), and \textbf{Poni} (manual work --- barber, cleaning).
\item It is \textbf{not a system of division of labour} (division of labour is economic and based on choice/freedom): it is based on \textbf{caste norms and religious significance}, and is what Ambedkar called \textbf{graded inequality} --- privileges for some castes, pathos for others --- in effect a \textbf{form of slavery for the Ponis}.
\item \textbf{Surinder Jodhka} found the jajmani relation \textbf{disintegrating} in Green Revolution areas: the \textbf{caste--occupation link is dissociating} (wage labour, monetised economy), with \textbf{marketisation of caste occupations} (Dalits opening barber shops for migrant labourers), landlords dropping attached labour for modern ways, and Dalit families doing domestic work in landlord homes --- all castes also \textbf{Sanskritising}.
\end{itemize}
\end{KeyConcept}

\begin{QuickRevision}
\begin{itemize}
\item Jajmani: Jajman, Purohit, Mahabrahmin, Poni; NOT division of labour; graded inequality (Ambedkar), slavery for Ponis.
\item Jodhka: jajmani disintegrating --- caste-occupation dissociation, marketisation, Sanskritisation.
\end{itemize}
\end{QuickRevision}

\section{Bonded Labour and New Forms of Bondage}

\begin{KeyConcept}
\begin{itemize}
\item Indian agriculture ran on \textbf{bonded (intergenerational) labour}. \textbf{Swami Agnivesh} fought for its removal (filing PILs in the Supreme Court; emancipating bonded labourers in Haryana, Chhattisgarh, Bihar); the practice was made illegal by law and court decisions. (The idea of 'bondage' now extends beyond agriculture --- e.g. an employer demanding a five-lakh bond is also illegal per the Supreme Court.)
\item \textbf{Attached labour} was also a form of bondage. After Green Revolution, \textbf{new forms of bondage} emerged: landowners gave (interest-free) loans to attach labourers --- 'work until you repay'; \textbf{Pranab Bardhan} found labourers' money incomes rose but \textbf{inflation cut their real purchasing power}; \textbf{Paul Brass} called the attached labourers \textbf{'bonded slaves'}.
\item \textbf{Jodhka's critique}: unlike permanent slavery, this bondage is \textbf{temporary} --- once the (interest-free) loan is repaid they are free; so it is better called a \textbf{'labour mortgage'} (the interest-free loan's purpose is to attach, not to profit).
\item In textile industries (Tamil Nadu, Nepal, Bangladesh) the same pattern appeared --- moneylenders (in collusion with mill owners) attached workers; when the WTO ruling removed the US quota, \textbf{Nepal's textile mills shut and the workers brought the Maoist revolution}.
\item A new bondage in GR areas: as migrant labourers declined, landowners \textbf{gave land and brought whole families} --- the family became servants of the landowner. \textbf{Beteille} also noted the rise of \textbf{formal contracts} between tenant and landowner.
\end{itemize}
\end{KeyConcept}

\begin{QuickRevision}
\begin{itemize}
\item Bonded labour (intergenerational; Swami Agnivesh's PILs); attached labour = bondage; new bondage via interest-free loans.
\item Bardhan: money income up but inflation cut real wages; Brass: 'bonded slaves'; Jodhka: temporary 'labour mortgage'.
\item Textile bondage (Tamil Nadu/Nepal/Bangladesh; WTO quota $\rightarrow$ Nepal's Maoist revolution); whole-family servitude; Beteille's formal contracts.
\end{itemize}
\end{QuickRevision}

\section{Agrarian Class Structure: Desai, Thorner and the Mode of Production Debate}

\begin{KeyConcept}
\begin{itemize}
\item \textbf{A. R. Desai}: Indian agriculture has \textbf{four classes} --- three cultivators (\textbf{landowners, tenants, agricultural labourers}) plus the \textbf{non-agriculturalist} (who services the agrarian class). These classes emerged with the \textbf{commercialisation of agriculture} under the British (permanent settlement; the Ryotwari system's exaggerated taxation $\rightarrow$ middlemen became big landlords $\rightarrow$ the Deccan riots; the Mahalwari system, where village elders kept the best land $\rightarrow$ a new class of \textbf{small/marginal farmers} and \textbf{big tenants} emerged).
\item \textbf{Daniel Thorner} rejected this --- one person can hold \textbf{multiple roles} (landowner can also be tenant; labourer can also be tenant) --- hence his famous \textbf{Malik, Kisan, Mazdoor}.
\item \textbf{Mode of production} (Marx) = the \textbf{means of production + relations of production} (and their \textbf{inherent contradictions}); human labour is both economic and social (the expression of human potential --- the dialectic of objectification). Marx held that at one time there is \textbf{one dominant mode of production}; the \textbf{Indian mode-of-production debate} asks whether, after Green Revolution, agriculture became \textbf{capitalist, or remained feudal, or is semi-feudal/semi-capitalist}.
\end{itemize}
\end{KeyConcept}

\begin{QuickRevision}
\begin{itemize}
\item A. R. Desai: 4 classes (landowner, tenant, labourer, non-agriculturalist); from commercialisation; Ryotwari $\rightarrow$ middlemen/big landlords (Deccan riots); Mahalwari $\rightarrow$ small/marginal + big tenants.
\item Thorner: multiple roles $\rightarrow$ Malik/Kisan/Mazdoor.
\item Mode of production = means + relations of production; one dominant mode (Marx); Indian debate: capitalist vs feudal vs semi-feudal.
\end{itemize}
\end{QuickRevision}
